% Quiz: Prompt Engineering (Weeks 1 & 2) - Multiple Choice (20 questions)
\documentclass[11pt,a4paper]{article}
\usepackage[utf8]{inputenc}
\usepackage{geometry}
\usepackage{enumitem}
\geometry{margin=1in}
\begin{document}
\begin{center}
  {\LARGE Prompt Engineering \textemdash{} Quiz 1}\\[6pt]
  {\large Name: \underline{\hspace{6cm}} \quad Date: \underline{\hspace{3cm}}}
\end{center}
\vspace{0.5cm}

\textbf{Instructions:} This quiz contains 20 multiple-choice questions. Answer all questions. Each question is worth 1 point unless otherwise noted. Total: 20 points.
\vspace{0.5cm}
\begin{enumerate}[leftmargin=*, label=\textbf{\#\arabic*}]

  \item Which of the following best describes a \textbf{prompt}?\\
    \begin{enumerate}[label=\Alph*)]
      \item A set of numerical parameters inside an LLM.
      \item The textual instruction, question, or example you give an AI system.
      \item The final output produced by an AI model.
      \item The hardware used to train a model.
    \end{enumerate}

  \item The simplest way to explain how most LLMs work is:\\
    \begin{enumerate}[label=\Alph*)]
      \item They store exact answers to every possible question.
      \item They predict the next token (word piece) based on patterns learned from large text corpora.
      \item They run symbolic logic rules coded by engineers.
      \item They rely only on live internet access to answer questions.
    \end{enumerate}

  \item Tokenization can split a single long word into several sub-word tokens.\\
    \begin{enumerate}[label=\Alph*)]
      \item True
      \item False
      \item Sometimes, but only for English
      \item Only for character-level tokenizers
    \end{enumerate}

  \item What does self-attention allow a Transformer model to do?\\
    \begin{enumerate}[label=\Alph*)]
      \item Compare and weigh other tokens so each token's representation becomes context-aware.
      \item Store fixed dictionary definitions for words.
      \item Translate text into image features.
      \item Ignore long-range dependencies.
    \end{enumerate}

  \item Which statement about embeddings is correct?\\
    \begin{enumerate}[label=\Alph*)]
      \item Embeddings are human-written definitions attached to each word.
      \item Embeddings are the final tokens produced by the model.
      \item Embeddings are numeric vectors that locate tokens in a ``meaning space'' so similar words are near each other.
      \item Embeddings are images generated by multimodal models.
    \end{enumerate}

  \item What does ``hallucination'' mean in the context of LLMs?\\
    \begin{enumerate}[label=\Alph*)]
      \item A visual artifact in image synthesis.
      \item The model producing plausible-sounding but incorrect or fabricated information.
      \item A model refusing to answer due to safety filters.
      \item When a model repeats the input prompt verbatim.
    \end{enumerate}

  \item Increasing model size (more parameters) always guarantees fewer errors and no hallucinations.\\
    \begin{enumerate}[label=\Alph*)]
      \item True
      \item False
      \item Only if trained on perfect data
      \item Only for language tasks
    \end{enumerate}

  \item Which of these is a useful prompt ``lever'' to improve AI output?\\
    \begin{enumerate}[label=\Alph*)]
      \item Role (e.g., ``You are a teacher'')
      \item Output format (e.g., ``Return 3 bullets'')
      \item Audience (e.g., ``for 9th graders'')
      \item All of the above
    \end{enumerate}

  \item In the self-attention mechanism each token computes Query, Key, and Value. Which description is most accurate?\\
    \begin{enumerate}[label=\Alph*)]
      \item Query = what I'm looking for; Key = what I contain; Value = the information I provide to others.
      \item Query = the token index; Key = the token frequency; Value = the token length.
      \item Query = model parameters; Key = training data; Value = learning rate.
      \item All three are identical and interchangeable.
    \end{enumerate}

  \item Which of the following best describes "fine-tuning"?\\
    \begin{enumerate}[label=\Alph*)]
      \item Training a model from scratch on random data.
      \item Additional training on curated examples to specialize or correct model behavior.
      \item Reducing the number of parameters to make the model smaller.
      \item Encrypting the model weights for security.
    \end{enumerate}

  \item What is an example of an emergent ability in large LLMs?\\
    \begin{enumerate}[label=\Alph*)]
      \item Decreased compute requirements.
      \item Suddenly solving tasks with few examples (few-shot learning) not seen in smaller models.
      \item The model physically improving hardware.
      \item Eliminating the need for tokenization.
    \end{enumerate}

  \item Why are position embeddings used in Transformers?\\
    \begin{enumerate}[label=\Alph*)]
      \item To encode the order of tokens since attention is permutation-invariant.
      \item To compress tokens into fewer dimensions.
      \item To convert text into images.
      \item To increase the model temperature.
    \end{enumerate}

  \item Which of these is NOT a typical limitation of LLMs?\\
    \begin{enumerate}[label=\Alph*)]
      \item Hallucination
      \item Biases inherited from training data
      \item Unlimited context-window (they can process arbitrarily long documents)
      \item Outdated knowledge due to training cutoff
    \end{enumerate}

  \item Tokenization helps models because: \\ 
    \begin{enumerate}[label=\Alph*)]
      \item It allows handling rare or long words by splitting them into familiar subparts.
      \item It prevents models from learning anything new.
      \item It ensures outputs are always grammatically correct.
      \item It stores complete books inside tokens.
    \end{enumerate}

  \item Which is a correct example of a prompt that uses constraints effectively?\\
    \begin{enumerate}[label=\Alph*)]
  \item ``Write about photosynthesis.''
  \item ``You are a high-school biology teacher. Explain photosynthesis in 3 bullets for 9th graders in under 100 words.''
  \item ``Do anything you want.''
  \item ``Make me rich.''
    \end{enumerate}

  \item Which metric describes how an embedding represents a token?\\
    \begin{enumerate}[label=\Alph*)]
      \item A high-dimensional numeric vector where nearby vectors indicate semantic similarity.
      \item A single integer representing word length.
      \item A boolean flag marking stopwords.
      \item A 2D image file.
    \end{enumerate}

  \item What role does ``instruction tuning'' play in model behavior?\\
    \begin{enumerate}[label=\Alph*)]
      \item It teaches the model to follow human instructions better via examples of prompts and desired outputs.
      \item It deletes the model's training data.
      \item It converts the model into a decision tree.
      \item It only changes the model's hardware.
    \end{enumerate}

  \item Which mitigation is effective against hallucinations?\\
    \begin{enumerate}[label=\Alph*)]
      \item Verifying facts with reliable external sources or adding a retrieval/citation step.
      \item Asking the model to guess more often.
      \item Increasing temperature to the maximum.
      \item Removing tokenization.
    \end{enumerate}

  \item What is the purpose of multi-head attention in Transformers?\\
    \begin{enumerate}[label=\Alph*)]
      \item It allows the model to attend to information from multiple representation subspaces and capture different relationships in parallel.
      \item It reduces the number of tokens by merging them.
      \item It encrypts the model weights for secure storage.
      \item It converts tokens into images for multimodal tasks.
    \end{enumerate}

  \item Which three elements do good prompts typically provide?\\
    \begin{enumerate}[label=\Alph*)]
      \item Purpose, context, and constraints.
      \item Tokens, parameters, and gradients.
      \item Hardware, dataset, and compute budget.
      \item Output format only.
    \end{enumerate}

\end{enumerate}

\vspace{0.5cm}
	\textbf{Scoring:} Each question = 1 point. Total points: 20.
\end{document}
